\begin{enumerate}
  \item
  \[
    f_e = \frac{f_1 f_2}{f_1 + f_2 - d} = \frac{20 \cdot (-10)}{20 - 10 - 15} = \frac{-200}{-5} = 40\,\mathrm{cm}.
  \]
  \item Lins 1 avbildar objektet i sitt brännplan, $20\,\mathrm{cm}$ bakom lins 1, dvs.\ $5\,\mathrm{cm}$ \emph{bakom} lins 2. Alltså är $s_2 = d - f_1 = -5\,\mathrm{cm}$ (ett virtuellt objekt). Linsformeln ger
  \[
    s_2' = \frac{s_2 f_2}{s_2 - f_2} = \frac{(-5)(-10)}{-5 + 10} = 10\,\mathrm{cm}.
  \]
  \item Totala längden är $15 + 10 = 25\,\mathrm{cm}$, medan en enkel lins med $f = 40\,\mathrm{cm}$ kräver $40\,\mathrm{cm}$ från linsen till sensorn.
  \item Månen är ett objekt i oändligheten, så $h' = f_e \tan\alpha$ med $\alpha = 0{,}26^\circ$ (från axeln till månens kant). Diametern blir
  \[
    2 f_e \tan 0{,}26^\circ = 2 \cdot 400\,\mathrm{mm} \cdot \tan 0{,}26^\circ \approx 3{,}6\,\mathrm{mm}.
  \]
  Teleobjektivet ger alltså samma bildstorlek som en enkel lins med $f = 40\,\mathrm{cm}$, men i ett betydligt kortare paket.
\end{enumerate}
